\section{Introduction}
\label{sec:intro}

Street-view perception models predict subjective attributes such as
safety, wealth, or liveliness at
scale~\cite{salesses2013,naik2014,dubey2016}, yet they remain limited as
explanatory tools: they do not answer the question central to urban
design, namely \emph{for this specific scene, which visual elements
drive the perception, and what targeted modification would plausibly
shift it?} Existing explainability methods such as saliency, SHAP, and
concept probes remain correlational, identifying features associated
with a score without validating that manipulating them changes human
judgement~\cite{adebayo2018,kim2018}. Empirical urban-design research
links greenery, maintenance, and visibility cues to perceived
safety~\cite{jacobs1961,newman1972,li2015greenery,jiang2018brokenwindows,portnov2020lighting},
but does not by itself establish which localized visual changes validly
shift perception for a specific scene.

Following Jacobs and defensible-space
theory~\cite{jacobs1961,newman1972}, urban perception is often discussed
as depending on multiple visible cues within the same scene. We test a
restricted version of this idea: whether a given scene admits multiple
distinct \emph{single-lever} interventions that each, evaluated
independently against the same unedited original, plausibly shift
perception.

We propose an \emph{interventional attribution} protocol that makes
this question empirically testable. Rather than treating arbitrary
regions as the explanatory unit, we define structured lever
interventions drawn from an editor-feasible intervention vocabulary,
each specifying a semantic concept, spatial support, intervention
direction, and constrained edit template. Candidate edits are realised
via prompt-conditioned image editing and retained only if a
vision--language critic judges them to preserve the same place, remain
localized, and appear realistic and plausible. Because generative
editors may introduce correlated non-target modifications, each edit is
treated as a hypothesis to be audited rather than as faithful evidence
by default.

This paper does not claim validated causal effects of specific edits on
human perception. Instead, it introduces an auditable interventional
attribution framework and a pilot benchmark for studying that question.
The current study uses a validity-gated generation pipeline and an
auxiliary perception proxy to prioritize promising interventions for
future human evaluation. Our empirical questions are: which lever
families produce directional improvement on an auxiliary proxy once
edits pass validity auditing, and across scenes, how many distinct
single-lever interventions remain promising after validity screening and
thresholding?
